% Chapter 14 of Genesis
\bookchapter{\name{Lot}'s Captivity and Rescue}

% verse 1
\bverse And it came to pass in the days of \name{Amraphel} king of \location{Shinar}, \name{Arioch} king of \location{Ellasar}, \name{Chedorlaomer} king of \location{Elam}, and \name{Tidal} king of nations, 
% verse 2
\bverse \that they made war with \name{Bera} king of \location{Sodom}, \name{Birsha} king of \location{Gomorrah}, \name{Shinab} king of \location{Admah}, \name{Shemeber} king of \location{Zeboiim}, and the king of \location{Bela} (that is, \name{Zoar}).
% verse 3
\bverse All these joined together in the \location{Valley of Siddim} (that is, the Salt Sea).
% verse 4
\bverse Twelve years they served \name{Chedorlaomer}, and in the thirteenth year they rebelled.
% verse 5
\bverse In the fourteenth year, \name{Chedorlaomer} and the kings that \were with him came and attacked the \people{Rephaim} in \location{Ashteroth Karnaim}, the \people{Zuzim} in \location{Ham}, the \people{Emim} in \location{Shaveh Kiriathaim},
% verse 6
\bverse and the \people{Horites} in their mountain of \location{Seir}, as far as El Paran, which \is by the wilderness.
% verse 7
\bverse Then they turned back and came to \location{En Mishpat} (that \is, \location{Kadesh}), and attacked all the country of the \people{Amalekites}, and also the \people{Amorites} who dwelt in \location{Hazezon Tamar}.

\todo{Try to make sense of all these people, places, peoples, and what is going on - perhaps try drawing some diagrams?}

% verse 8
\bverse And the king of \location{Sodom}, the king of \location{Gomorrah}, the king of \location{Admah}, the kings of Zeboiim, and the king of \location{Bela} (that \is, \name{Zoar}) went out and joined together in battle in the \location{Valley of Siddim}
% verse 9
\bverse against \name{Chedorlaomer} king of \location{Elam}, \name{Tidal} king of nations, \name{Amraphel} king of \location{Shinar}, and \name{Arioch} king of \location{Ellasar}--four kings against five.
% verse 10
\bverse Now the \location{Valley of Siddim} \textit{was full of} asphalt pits; and the kings of \location{Sodom} and \location{Gomorrah} fled; \textit{some} fell there, and the remainder fled to the mountains.
% verse 11
\bverse Then they took all the goods of \location{Sodom} and \location{Gomorrah}, and all their provisions, and went their way.
% verse 12
\bverse They also took \name{Lot}, \name{Abram};s brother's son who dwelt in \location{Sodom}, and his goods, and departed.
% verse 13
\bverse Then one who had escaped came and told \name{Abram} the Hebrew, for he dwelt by the terebinth trees of \name{Mamre} the \people{Amorite}, brother of \name{Eshcol} and brother of \name{Aner}; and they \were allies with \name{Abram}.
% verse 14
\bverse Now when \name{Abram} heard that his brother was taken captive, he armed his three hundred and eighteen trained \textit{servants} who were born in his own house, and went in pursuit as far as \location{Dan}.
% verse 15
\bverse He divided his forces against them by night, and he and his servants attacked them and pursued them as far as \location{Hobah}, which \is north of \location{Damascus}.
% verse 16
\bverse So he brought back all the goods, and also brought back his brother \name{Lot} and his goods, as well as the women and the people.
% verse 17
\bverse And the king of \location{Sodom} went out to meet him at the Valley of Shaveh (that \is, the King's Valley), after his return from the defeat of \name{Chedorlaomer} and the kings who \were with him.

\bmarkerdown{\name{Abram} and \name{Melchizedek}}

% verse 18
\bverse Then \name{Melchizedek} king of \location{Salem} brought out bread and wine; he \was the priest of God Most High\vmark{a}.
	\translationnote{a}{This is the first time the term `Most High' has appeared. This term comes from the Hebrew term `elyown' (Strong's 5945 \cite{Hebrew Concordance}), which is an elevating title meaning \textit{the Supreme} or the \textit{Uppermost}. This is a differentiating statement which means this is referring to the literal \textit{highest} being above all others.}
	\doctrinenote{a}{This is a supporting concept to their being multiple members of the `God' family, or multiple `God' beings. This aligns with the plurality used in \bref{Genesis 1:26}, \bref{John 1:1}, \bref{John 17:5} and other places and perhaps points to the idea of a hierarchy in the God family.}

% verse 19
\bverse And he blessed him and said:
\begin{bquotation}
% verse 20
``Blessed be \name{Abram} of God Most High, Possessor of heaven and earth; \bverse And blessed be God Most High, Who has delivered your enemies into your hand.''
\end{bquotation}
And he gave him a tithe of all.
% verse 21
\bverse Now the king of \location{Sodom} said to \name{Abram}, ``Give me the persons, and take the goods for yourself.''
% verse 22
\bverse But \name{Abram} said to the king of \location{Sodom}, ``I have raised my hand to the \lord, God Most High, the Possessor of heaven and earth,
% verse 23
\bverse that I \textit{will take} nothing, from a thread to a sandal strap, and that I will not take anything that \is yours, lest you should say, `I have made \name{Abram} rich'--
% verse 24
\bverse except only what the young men have eaten, and the portion of the man who went with me: \name{Aner}, \name{Eshcol}, and \name{Mamre}; let them take their portion.''
